\section*{\centering LỜI MỞ ĐẦU}
\addcontentsline{toc}{section}{Lời mở đầu}

Hoạt động nghiên cứu khoa học ở trường đại học gắn với nhiều công việc, từ quản lý hồ sơ nhà nghiên cứu đến theo dõi đề tài, phân bổ kinh phí và đánh giá kết quả. Mỗi đề tài có thể liên quan đến nhiều người, nhiều đơn vị và nhiều nguồn tài trợ. Thông tin phát sinh ở các giai đoạn khác nhau cần được tập hợp để người phụ trách có thể theo dõi quá trình thực hiện và xem xét kết quả khi nghiệm thu.

Nhóm chọn đề tài \textit{\ReportTitle} để vận dụng kiến thức trong học phần Lập trình hướng đối tượng. Nhà nghiên cứu, đề tài và hội đồng nghiệm thu có thông tin, trách nhiệm và quan hệ riêng. Những đặc điểm đó tạo điều kiện để nhóm thực hành thiết kế lớp, đóng gói, trừu tượng hóa, kế thừa và đa hình trong cùng một bài toán.

Trong phạm vi bài tập lớn, nhóm xây dựng chương trình bằng C\# để mô phỏng các nghiệp vụ về nhân sự, đề tài, kinh phí, hội đồng nghiệm thu và kết quả nghiên cứu. Khi triển khai, nhóm chú trọng phân chia trách nhiệm giữa các lớp và lựa chọn mẫu thiết kế phù hợp với từng nghiệp vụ.

Báo cáo trình bày các yêu cầu của bài toán, những lựa chọn trong thiết kế và cài đặt, cùng kết quả kiểm thử chương trình.
